\DocumentMetadata{
  pdfversion=2.0,
  pdfstandard=ua-2,
  lang=en-US,
  tagging=on,
  tagging-setup={math/setup=mathml-SE,tabsorder=structure}
}
\documentclass[11pt,letterpaper]{article}
\usepackage[margin=0.68in,includefoot]{geometry}
\usepackage{newcomputermodern}
\usepackage{amsmath,amscd,mathtools}
\usepackage{tikz,adjustbox}
\providecommand{\square}{\mdlgwhtsquare}
\usepackage[hidelinks]{hyperref}
\hypersetup{
  pdftitle={Algebra PhD Exam, July 1997},
  pdfauthor={University of Florida Department of Mathematics},
  pdfsubject={Graduate mathematics examination},
  pdfdisplaydoctitle=true,
  bookmarksopen=true
}
\setcounter{secnumdepth}{0}
\setlength{\parindent}{0pt}
\setlength{\parskip}{0.28em}
\setlength{\emergencystretch}{3em}
\setlength{\leftmargini}{2.15em}
\setlength{\leftmarginii}{2.65em}
\setlength{\leftmarginiii}{3em}
\pagestyle{empty}
\raggedbottom
\allowdisplaybreaks
\NewTaggingSocketPlug{block/list/label}{examproblem}{%
  \tagstructbegin{tag=Lbl}\tagstructend%
  \tagstructbegin{tag=\LBody}%
  \tagstructbegin{tag=\ProblemHeadingTag,title-o={\ProblemHeadingText}}%
  \tagmcbegin{tag=\ProblemHeadingTag,actualtext={\ProblemHeadingText}}%
  #2%
  \tagmcend%
  \tagstructend%
}
\newenvironment{examproblems}[2]{%
  \def\ProblemHeadingTag{#2}%
  \AssignTaggingSocketPlug{block/list/label}{examproblem}%
  \begin{enumerate}%
  \setcounter{enumi}{#1}%
  \renewcommand{\labelenumi}{\textbf{\theenumi.}}%
  \setlength{\topsep}{0.35em}%
  \setlength{\partopsep}{0pt}%
  \setlength{\itemsep}{0.48em}%
  \setlength{\parsep}{0.18em}%
}{\end{enumerate}}
\newenvironment{examsubparts}{%
  \AssignTaggingSocketPlug{block/list/label}{default}%
  \begin{enumerate}%
  \setlength{\topsep}{0.18em}%
  \setlength{\partopsep}{0pt}%
  \setlength{\itemsep}{0.16em}%
  \setlength{\parsep}{0.1em}%
}{\end{enumerate}}
\newenvironment{examinstructions}{%
  \par\begingroup\itshape
}{%
  \par\endgroup\vspace{0.18em}
}
\makeatletter
\renewcommand\subsection{\@startsection{subsection}{2}{0pt}%
  {-1.25ex plus -.25ex minus -.1ex}%
  {0.55ex plus .1ex}%
  {\normalfont\large\bfseries}}
\renewcommand{\@maketitle}{%
  \newpage\null\vskip -1em%
  {\footnotesize\bfseries
    \href{https://www.ufl.edu/}{University of Florida}, \href{https://math.ufl.edu/}{Department of Mathematics}\par}%
  \vskip -0.5em%
  \tagstructbegin{tag=H1,title={Algebra PhD Exam, July 1997}}%
  \tagpdfparaOff%
  \tagmcbegin{tag=H1}%
    {\LARGE\bfseries\href{https://gma.math.ufl.edu/exams/algebra-phd/}{Algebra PhD Exam}, July 1997\par}%
  \tagmcend%
  \tagpdfparaOn%
  \tagstructend%
  \vskip 0.25em%
  \tagmcbegin{artifact}%
    \hrule height 1pt%
  \tagmcend%
  \vskip 1em%
}
\makeatother
\title{Algebra PhD Exam, July 1997}
\author{}
\date{}
\begin{document}
\maketitle
\thispagestyle{empty}
\begin{examinstructions}
Answer seven out of the following ten problems. State clearly any results you need.
\end{examinstructions}
\begin{examproblems}{0}{H2}
\def\ProblemHeadingText{Problem 1.}
\item
Let~\(R\) be a commutative ring with identity, and~\(S\) a multiplicative system in~\(R\). Show that the localization \(S^{-1}R\) is a flat~\(R\)-module.
\def\ProblemHeadingText{Problem 2.}
\item
Show that the field of complex numbers has uncountably many distinct automorphisms.
\def\ProblemHeadingText{Problem 3.}
\item
This problem has two parts.
\begin{examsubparts}
\item[\textnormal{(a)}]
Show that \(A_5\) is simple.
\item[\textnormal{(b)}]
Use (a) to show that \(S_n\) is not solvable for \(n \geq 5\).
\end{examsubparts}
\def\ProblemHeadingText{Problem 4.}
\item
If~\(k\) is a field, show that the ring of power series \(k[[T]]\) is a discrete valuation ring.
\def\ProblemHeadingText{Problem 5.}
\item
Suppose that~\(p\) and~\(q\) are distinct primes such that \(q>p^2\) and \(q \not\equiv 1 \pmod p\). Show that if~\(G\) is a group of order \(p^2q\), then~\(G\) is abelian.
\def\ProblemHeadingText{Problem 6.}
\item
Let~\(R\) be a local Noetherian (commutative) ring, and let~\(M\) be a finitely generated~\(R\)-module. Show that~\(M\) is flat if and only if it is free.
\def\ProblemHeadingText{Problem 7.}
\item
Suppose that~\(R\) is a ring with the property that every short exact sequence of unital~\(R\)-modules splits. Show that every unital~\(R\)-module is a direct sum of simple unital~\(R\)-modules.
\def\ProblemHeadingText{Problem 8.}
\item
This problem has two parts.
\begin{examsubparts}
\item[\textnormal{(a)}]
Show that if~\(k\) is a perfect field, then an irreducible polynomial in \(k[T]\) has no multiple roots. (Hint: show that if \(f(T)\) is irreducible and has multiple roots, then the characteristic of~\(k\) is \(p>0\), and \(f(T)=g(T^p)\) for some \(g\in k[T]\).)
\item[\textnormal{(b)}]
Give an example of a (non-perfect) field~\(k\) and an irreducible polynomial \(f(T)\) with multiple roots.
\end{examsubparts}
\def\ProblemHeadingText{Problem 9.}
\item
Suppose~\(R\) is a ring with identity. For each left~\(R\)-module~\(M\), show that \(\operatorname{Hom}_R(R,M)\) is naturally isomorphic to~\(M\) (in the process, explain what “naturally isomorphic” means).
\def\ProblemHeadingText{Problem 10.}
\item
Classify up to isomorphism all semisimple rings of order \(1584\).
\end{examproblems}
\end{document}
